\documentclass[11pt]{article}

\usepackage[utf8]{inputenc}
\usepackage[T1]{fontenc}
\usepackage{lmodern}
\usepackage{geometry}
\usepackage{amsmath,amssymb,amsfonts,amsthm,mathtools}
\usepackage{bm}
\usepackage{enumitem}
\usepackage{microtype}
\usepackage{hyperref}
\usepackage{titlesec}
\usepackage{booktabs}
\usepackage{array}
\usepackage{setspace}
\usepackage{mathrsfs}
\usepackage{physics}

\geometry{margin=1in}
\hypersetup{colorlinks=true, linkcolor=black, urlcolor=blue, citecolor=black}
\setlist[enumerate]{topsep=0.4em,itemsep=0.35em}
\setlist[itemize]{topsep=0.3em,itemsep=0.25em}
\titleformat{\section}{\large\bfseries}{\thesection.}{0.5em}{}
\titleformat{\subsection}{\normalsize\bfseries}{\thesubsection.}{0.5em}{}
\setstretch{1.06}

\newcommand{\Aether}{\AE{}ther}
\newcommand{\Aflow}{\AE{}ther-flow}
\newcommand{\TheoryName}{\Aflow{} Interpretation of Relativity}
\newcommand{\TheoryProgram}{\Aether{} / \Aflow{} framework}
\newcommand{\Sub}{\mathcal{A}}
\newcommand{\BridgeMetric}{\mathfrak{g}}
\newcommand{\Ell}{\mathcal{L}}

\newtheorem{definition}{Definition}
\newtheorem{proposition}{Proposition}
\newtheorem{remark}{Remark}
\newtheorem{corollary}{Corollary}
\newtheorem{theorem}{Theorem}

\title{\textbf{The \TheoryName{}}\\[0.4em]
\large Higher-Derivative Quartic Branch Explicit Bridge-Tensor Reduction Test}
\author{}
\date{}

\begin{document}

\maketitle

\begin{abstract}
The previous bridge-closure note isolated the exact remaining derivational burden on the surviving quartic branch: compute the observer-equation correction tensor
\[
G_{\mu\nu}[g]
=
8\pi G_N\,T_{\mu\nu}^{\mathrm{matter}}
+ \mathcal C_{\mu\nu}^{\mathrm{br}}
\]
for the explicit fixed completion \(S_{\mathrm{min},4}\), and decide whether that tensor vanishes, is only gauge/constraint exact, or leaves a genuine observer-accessible remainder. This note performs that next step at the present disciplined scope by combining the explicit-completion reduced-system note, the unit-lapse weighted/homogeneous same-class closure line, and the quartic recovery-compatibility reduction.

The resulting answer is negative for exact derivational closure. After eliminating the \(Q\)-sector and ordered-flow auxiliaries through their reduced equations, the surviving branch correction tensor is exactly the effective quartic-scalar stress tensor:
\[
\mathcal C_{\mu\nu}^{\mathrm{br}}
=
8\pi G_N\,T_{\mu\nu}^{(\sigma,\mathrm{eff})}.
\]
In reduced \(3+1\) variables its leading components are
\begin{align*}
\rho_\sigma
&=
\frac{1}{2m_S^2}\pi_\sigma^2
+ \frac{\lambda_4}{2M_*^2}(\Delta_\gamma\sigma)^2
+ \mathcal O_{\mathrm{l.o.t.}},
\\
J_i^{(\sigma)}
&=
-\pi_\sigma D_i\sigma
+ \mathcal O_{\mathrm{l.o.t.}},
\\
S_{ij}^{(\sigma)}
&=
\frac{\lambda_4}{M_*^2}
\left(
D_iD_j\sigma\,\Delta_\gamma\sigma
-\frac{1}{2}\gamma_{ij}(\Delta_\gamma\sigma)^2
\right)
+ \mathcal O_{\mathrm{l.o.t.}}.
\end{align*}
On the decaying flat branch, eliminating the remaining quartic scalar mode yields the explicit scalar self-energy
\[
\Delta\mathcal L_{\phi,4}^{(2)}
=
-\frac{1}{2}\Pi_\phi^{(4)}(\omega,\mathbf{k})\,|\phi|^2,
\qquad
\Pi_\phi^{(4)}
=
m_S^2\sigma_0^2
\frac{(\lambda_4/M_*^2)k^4}
{m_S^2\omega^2+(\lambda_4/M_*^2)k^4}.
\]
Because this scalar projection is nonzero for generic \((\omega,\mathbf{k})\) on the healthy quartic branch, the bridge correction tensor neither vanishes nor reduces to a pure gauge/constraint-exact term. The surviving quartic branch therefore remains a stable recovery-compatible candidate completion, but it does not close onto the exact Einsteinian observer equation at the present disciplined scope. In that precise sense, the first genuine derivational obstruction appears here, not at the same-class nonlinear-stability stage.
\end{abstract}

\section{Introduction}

The active manuscript sequence of \TheoryName{} has now reached a much sharper endpoint on the surviving quartic branch than before. The same fixed completion \(S_{\mathrm{min},4}\) has a reduced observer-accessible \(3+1\) system, a first local theorem, a coercive bootstrap, a repaired weighted/homogeneous unit-lapse route, and a positive same-class asymptotic/nonlinear-stability closure result. The previous bridge-closure note then isolated the exact remaining burden: compute the observer-side bridge correction tensor itself rather than asking for another stability repair.

This note performs that calculation at the present disciplined scope. The question is no longer whether the surviving branch is dynamically viable in the same class. It is whether, after the auxiliary sectors are genuinely reduced, the observer equation becomes exactly Einsteinian or whether a residual metric-mediated quartic correction survives.

\paragraph{Framework and claim boundary.}
\TheoryProgram{} denotes the broader \Aether{} / \Aflow{} framework built on the claims that the \Aether{} is the underlying four-dimensional substrate of reality, the \Aflow{} is its intrinsic ordered motion, observed three-dimensional space is the local experiential slice of that deeper substrate, S-time is the experienced order of change arising from matter, light, and the \Aflow{}, and observed expansion is the three-dimensional appearance of deeper four-dimensional ordered motion. Within that framework, \TheoryName{} denotes the exact-closure theory in which the effective gravitational dynamics are adopted to be exactly Einsteinian with universal matter coupling. In the active sequence, \emph{adoption} means use of that established relativistic dynamics without claiming substrate derivation, while \emph{derivation} is reserved for a first-principles recovery from explicit substrate variables.

The present note stays strictly within that claim boundary. It does not claim a completed substrate derivation of Einsteinian gravity. It does carry the recorded bridge-closure program through the next exact test: explicit reduction of the branch correction tensor \(\mathcal C_{\mu\nu}^{\mathrm{br}}\) on the surviving quartic branch and an explicit verdict on its status.

\section{Current Branch Data and the Reduction Target}

\begin{definition}[Surviving quartic branch data]
\label{def:branch}
The \emph{surviving quartic branch data} consist of:
\begin{enumerate}
    \item the bridge metric
    \begin{equation}
    \BridgeMetric_{AB}=G_{AB}-2U_AU_B,
    \label{eq:bridge}
    \end{equation}
    and observer metric
    \begin{equation}
    g_{\mu\nu}=\Xi^\ast\BridgeMetric_{AB},
    \label{eq:obsmetric}
    \end{equation}
    \item the universal matter route
    \begin{equation}
    S_{\mathrm{matter}}[\Xi^\ast\BridgeMetric,\psi]
    =
    S_{\mathrm{matter}}[g,\psi],
    \label{eq:matterroute}
    \end{equation}
    \item the explicit minimal quartic-compatible completion \(S_{\mathrm{min},4}\),
    \begin{equation}
    S_{\mathrm{min},4}
    =
    \frac{1}{16\pi G_N}
    \int_{\Sub} d^4X\,\sqrt{G}\,R[\BridgeMetric]
    + \int_{\Sub} d^4X\,\sqrt{G}\,\mathcal L_{\mathrm{aux,min},4}
    + S_{\mathrm{matter}}[g,\psi],
    \label{eq:Smin4}
    \end{equation}
    with
    \begin{equation}
    \mathcal L_{\mathrm{aux,min},4}
    \supset
    -\frac{m_S^2}{2}(U^A\nabla_A S-\sigma_0)^2
    -\frac{\lambda_4}{2M_*^2}(\Delta_\perp S)^2,
    \label{eq:quarticpiece}
    \end{equation}
    together with positive \(Q\)-sector and ordered-flow auxiliary blocks,
    \item the fixed-completion unit-lapse weighted/homogeneous same-class asymptotic closure route.
\end{enumerate}
\end{definition}

\begin{definition}[Bridge correction tensor]
\label{def:C}
For the branch data of Definition \ref{def:branch}, define \(\mathcal C_{\mu\nu}^{\mathrm{br}}\) by
\begin{equation}
G_{\mu\nu}[g]
=
8\pi G_N\,T_{\mu\nu}^{\mathrm{matter}}
+ \mathcal C_{\mu\nu}^{\mathrm{br}}.
\label{eq:bridgeeq}
\end{equation}
\end{definition}

\begin{remark}
The previous bridge-closure note showed that the exact derivational question is whether \(\mathcal C_{\mu\nu}^{\mathrm{br}}\) vanishes or is only gauge/constraint exact. The present note therefore has only one real task: reduce \(\mathcal C_{\mu\nu}^{\mathrm{br}}\) explicitly enough to settle that verdict.
\end{remark}

\section{Reduced Observer Equation After Auxiliary Elimination}

The explicit-completion reduced-system note already determines how the nonmetric sectors enter the observer-accessible equations. The point here is to reorganize that result directly at the level of \(\mathcal C_{\mu\nu}^{\mathrm{br}}\).

\begin{proposition}[Auxiliary elimination leaves only the quartic scalar channel]
\label{prop:auxelim}
On the surviving quartic branch, after imposing the reduced \(Q\)-sector and ordered-flow auxiliary equations, the bridge correction tensor decomposes as
\begin{equation}
\mathcal C_{\mu\nu}^{\mathrm{br}}
=
8\pi G_N\,T_{\mu\nu}^{(\sigma,\mathrm{eff})}
+ \mathcal E_{\mu\nu}^{(Q,\chi,W)},
\label{eq:Cdecomp0}
\end{equation}
where \(\mathcal E_{\mu\nu}^{(Q,\chi,W)}\) is a linear combination of the eliminated auxiliary equations and therefore vanishes on reduced solutions. Consequently,
\begin{equation}
\mathcal C_{\mu\nu}^{\mathrm{br}}
=
8\pi G_N\,T_{\mu\nu}^{(\sigma,\mathrm{eff})}
\qquad
\text{on the reduced surviving branch.}
\label{eq:CeqTeff}
\end{equation}
\end{proposition}

\begin{proof}
The explicit-completion reduced-system note shows that the \(Q\)-sector obeys an elliptic constraint with no linear flat-branch source and that the ordered-flow variables \((\chi,W_i^T)\) are likewise slice-wise auxiliary unknowns rather than observer-accessible propagating channels. In the flat reduction these sectors satisfy
\[
q^I=0,
\qquad
\nabla^2\chi=-\frac{1}{2}\partial_0H,
\qquad
V_i^T=-S_i,
\]
and leave no independent tensor or transverse-vector observer correction. More generally, varying the reduced action after substitution of auxiliary solutions produces the direct metric variation of the reduced quartic scalar action plus terms proportional to the auxiliary Euler--Lagrange equations. Those latter terms vanish on reduced solutions. Hence the only surviving observer-side branch correction is the effective quartic scalar stress tensor, which gives \eqref{eq:Cdecomp0} and \eqref{eq:CeqTeff}.
\end{proof}

\begin{corollary}[Reduced observer equation]
\label{cor:redeq}
The reduced observer equation on the surviving quartic branch is
\begin{equation}
G_{\mu\nu}[g]
=
8\pi G_N\left(
T_{\mu\nu}^{\mathrm{matter}}
+ T_{\mu\nu}^{(\sigma,\mathrm{eff})}
\right).
\label{eq:redeq}
\end{equation}
\end{corollary}

\begin{proof}
This is exactly \eqref{eq:bridgeeq} together with \eqref{eq:CeqTeff}.
\end{proof}

\begin{remark}
Corollary \ref{cor:redeq} is already the first serious verdict. The bridge correction tensor does not disappear automatically under auxiliary reduction. It becomes the effective quartic scalar stress tensor. The remaining question is therefore whether that stress tensor itself vanishes or is only gauge/constraint exact on the surviving branch.
\end{remark}

\section{Explicit Reduced Form of the Bridge Tensor}

Write the observer metric in \(3+1\) form,
\begin{equation}
g_{\mu\nu}dx^\mu dx^\nu
=
-N^2dt^2+\gamma_{ij}(dx^i+\beta^i dt)(dx^j+\beta^j dt),
\label{eq:ADM}
\end{equation}
with unit future normal \(n^\mu\), and decompose
\begin{equation}
S=\sigma_0 t+\sigma,
\qquad
\pi_\sigma
\equiv
m_S^2\bigl(n^\mu\nabla_\mu S-\sigma_0\bigr).
\label{eq:pisigma}
\end{equation}

\begin{proposition}[Explicit \(3+1\) decomposition of the bridge tensor]
\label{prop:3plus1}
On the reduced surviving quartic branch,
\begin{equation}
\mathcal C_{\mu\nu}^{\mathrm{br}}
=
8\pi G_N
\left(
\rho_\sigma n_\mu n_\nu
+2n_{(\mu}J_{\nu)}^{(\sigma)}
+S_{\mu\nu}^{(\sigma)}
\right),
\label{eq:C3plus1}
\end{equation}
where \(S_{\mu\nu}^{(\sigma)}n^\nu=0\), and the leading pieces are
\begin{align}
\rho_\sigma
&=
\frac{1}{2m_S^2}\pi_\sigma^2
+\frac{\lambda_4}{2M_*^2}(\Delta_\gamma\sigma)^2
+\mathcal O_{\mathrm{l.o.t.}},
\label{eq:rhosigma}
\\
J_i^{(\sigma)}
&=
-\pi_\sigma D_i\sigma
+\mathcal O_{\mathrm{l.o.t.}},
\label{eq:Jsigma}
\\
S_{ij}^{(\sigma)}
&=
\frac{\lambda_4}{M_*^2}
\left(
D_iD_j\sigma\,\Delta_\gamma\sigma
-\frac{1}{2}\gamma_{ij}(\Delta_\gamma\sigma)^2
\right)
+\mathcal O_{\mathrm{l.o.t.}}.
\label{eq:Sijsigma}
\end{align}
\end{proposition}

\begin{proof}
The explicit-completion reduced-system note already writes the Einstein constraints and evolution equations with the quartic scalar source separated from the ordinary matter source:
\begin{align*}
R[\gamma]+K^2-K_{ij}K^{ij}
&=
16\pi G_N\bigl(\rho^{\mathrm{matter}}+\rho_\sigma\bigr),
\\
D^j(K_{ij}-\gamma_{ij}K)
&=
8\pi G_N\bigl(J_i^{\mathrm{matter}}+J_i^{(\sigma)}\bigr),
\\
(\partial_t-\Ell_\beta)K_{ij}
&=
\cdots
-8\pi G_N
\left(
S_{ij}^{\mathrm{matter}}
-\frac{1}{2}\gamma_{ij}\bigl(S^{\mathrm{matter}}-\rho^{\mathrm{matter}}\bigr)
+\Theta_{ij}^{(\sigma)}
\right).
\end{align*}
Reassembling the normal--tangential decomposition of the Einstein equation gives \eqref{eq:C3plus1}. The formulas \eqref{eq:rhosigma} and \eqref{eq:Jsigma} are the recorded leading scalar source terms. The quartic spatial stress has the principal variation shown in \eqref{eq:Sijsigma}, while the lower-order commutator, curvature, and gauge-coupling terms are exactly the \(\mathcal O_{\mathrm{l.o.t.}}\) terms already grouped below the principal metric--scalar hierarchy in the reduced-system note. This yields \eqref{eq:C3plus1}--\eqref{eq:Sijsigma}.
\end{proof}

\begin{remark}
Proposition \ref{prop:3plus1} is the explicit observer-side tensor reduction needed for the present test. It shows directly that the bridge correction tensor carries the energy density, momentum density, and spatial stress of the quartic scalar channel. Unless that channel itself disappears, \(\mathcal C_{\mu\nu}^{\mathrm{br}}\) cannot vanish.
\end{remark}

\section{Flat-Branch Scalar Projection of the Bridge Tensor}

To decide whether the tensor in Proposition \ref{prop:3plus1} might still be pure gauge/constraint exact, one should look at the fully reduced flat-branch scalar projection where the auxiliary constraints are already solved.

\begin{proposition}[Flat-branch scalar projection]
\label{prop:scalarproj}
On the decaying flat branch with \(h_{00}=-2\phi\), auxiliary elimination yields the reduced scalar channel
\begin{equation}
\mathcal L_{\sigma,\mathrm{eff},4}^{(2)}
=
-\frac{m_S^2}{2}\bigl(\partial_0\sigma-\sigma_0\phi\bigr)^2
-\frac{1}{2}\sigma\,\mathcal B_4(-\nabla^2)\,\sigma,
\label{eq:Lsigmaeff}
\end{equation}
with
\begin{equation}
\mathcal B_4(-\nabla^2)
=
\frac{\lambda_4}{M_*^2}(-\nabla^2)^2.
\label{eq:B4op}
\end{equation}
Eliminating \(\sigma\) gives the explicit scalar bridge self-energy
\begin{equation}
\Delta\mathcal L_{\phi,4}^{(2)}(\omega,\mathbf{k})
=
-\frac{1}{2}\Pi_\phi^{(4)}(\omega,\mathbf{k})\,|\phi(\omega,\mathbf{k})|^2,
\label{eq:Deltaphi}
\end{equation}
where
\begin{equation}
\Pi_\phi^{(4)}(\omega,\mathbf{k})
=
m_S^2\sigma_0^2
\frac{\mathcal B_4(k^2)}
{m_S^2\omega^2+\mathcal B_4(k^2)}
=
m_S^2\sigma_0^2
\frac{(\lambda_4/M_*^2)k^4}
{m_S^2\omega^2+(\lambda_4/M_*^2)k^4}.
\label{eq:Piphi}
\end{equation}
Equivalently, the scalar projection of the linearized bridge tensor is
\begin{equation}
\delta\mathcal C_{\phi}^{\mathrm{br}}
=
\Pi_\phi^{(4)}(\omega,\mathbf{k})\,\phi(\omega,\mathbf{k}).
\label{eq:Cphi}
\end{equation}
\end{proposition}

\begin{proof}
Equations \eqref{eq:Lsigmaeff}--\eqref{eq:Piphi} are exactly the recorded quartic flat-branch reduction from the higher-derivative consistency and recovery-compatibility notes. The only surviving flat-background observer correction after auxiliary elimination is the scalar self-energy \(\Pi_\phi^{(4)}\). Varying \eqref{eq:Deltaphi} with respect to \(\phi\) yields the scalar projection \eqref{eq:Cphi} of the linearized bridge correction tensor.
\end{proof}

\begin{remark}
Equation \eqref{eq:Cphi} is the decisive observer-side invariant of the current note. It is already after auxiliary elimination, already on the reduced flat branch, and already in the scalar channel that survived all earlier coefficient-level tests. If this quantity is nonzero, then the bridge correction tensor carries genuine observer-side content.
\end{remark}

\section{Verdict on Vanishing and Gauge/Constraint Exactness}

\begin{definition}[Healthy nontrivial quartic regime]
\label{def:healthy}
The \emph{healthy nontrivial quartic regime} is defined by
\begin{equation}
m_S^2>0,
\qquad
\lambda_4>0,
\qquad
\sigma_0\neq 0,
\label{eq:healthy}
\end{equation}
together with the positive \(Q\)-sector mass block and the long-wavelength, non-ultrasoft recovery-compatible regime already recorded in the recovery-compatibility note.
\end{definition}

\begin{theorem}[Explicit bridge-tensor verdict]
\label{thm:verdict}
Assume the surviving quartic branch data of Definition \ref{def:branch} in the healthy nontrivial quartic regime of Definition \ref{def:healthy}. Then:
\begin{enumerate}
    \item \(\mathcal C_{\mu\nu}^{\mathrm{br}}\not\equiv 0\) on the surviving quartic branch.
    \item \(\mathcal C_{\mu\nu}^{\mathrm{br}}\) is not pure gauge/constraint exact.
    \item The surviving quartic branch therefore does not satisfy exact derivational bridge closure at the present disciplined scope.
\end{enumerate}
More precisely, the branch leaves a genuine observer-accessible remainder whose scalar projection is the quartic self-energy \(\Pi_\phi^{(4)}\), which is infrared-suppressed in the recorded recovery-compatible regime but not identically zero.
\end{theorem}

\begin{proof}
By Proposition \ref{prop:scalarproj},
\[
\delta\mathcal C_{\phi}^{\mathrm{br}}
=
\Pi_\phi^{(4)}(\omega,\mathbf{k})\,\phi(\omega,\mathbf{k}).
\]
In the healthy nontrivial quartic regime one has \(m_S^2>0\), \(\lambda_4>0\), and \(\sigma_0\neq 0\). Therefore for generic \(k\neq 0\),
\[
\Pi_\phi^{(4)}(\omega,\mathbf{k})>0.
\]
Hence the scalar projection of the bridge correction tensor is nonzero, so \(\mathcal C_{\mu\nu}^{\mathrm{br}}\) does not vanish identically. This proves item (1).

For item (2), suppose by contradiction that \(\mathcal C_{\mu\nu}^{\mathrm{br}}\) were pure gauge/constraint exact. Then on the decaying flat branch, after the auxiliary constraints are already solved and after fixing the ordinary scalar metric gauge used to define \(\phi\), the reduced scalar self-energy would have to vanish: pure gauge terms disappear after gauge fixing, and pure constraint terms disappear once the reduced constraints have been imposed. But Proposition \ref{prop:scalarproj} shows that the surviving scalar projection is precisely \(\Pi_\phi^{(4)}\phi\), which is nonzero for generic \((\omega,\mathbf{k})\). This contradiction proves that \(\mathcal C_{\mu\nu}^{\mathrm{br}}\) is not pure gauge/constraint exact.

Item (3) is then immediate from the bridge-closure criterion established in the previous bridge-closure note: exact derivational closure requires vanishing or pure gauge/constraint exactness, and neither holds here.
\end{proof}

\begin{corollary}[Current scientific status of the surviving quartic branch]
\label{cor:status}
At the present disciplined scope, the surviving quartic branch remains:
\begin{enumerate}
    \item \textbf{positive as a stability result:} the same-class weighted/homogeneous unit-lapse route closes globally;
    \item \textbf{positive as a recovery-compatibility result:} the residual observer correction is quartically suppressed in the long-wavelength, non-ultrasoft regime;
    \item \textbf{negative as an exact derivation result:} the bridge correction tensor leaves a genuine observer-accessible remainder and therefore does not close back onto the exact Einsteinian observer equation.
\end{enumerate}
\end{corollary}

\begin{proof}
The first statement is the recorded asymptotic/nonlinear-stability closure result. The second is the recovery-compatibility result. The third is Theorem \ref{thm:verdict}.
\end{proof}

\begin{corollary}[Location of the first genuine derivational obstruction]
\label{cor:firstobstruction}
The first genuine obstruction on the surviving quartic branch does not appear at the same-class nonlinear-stability stage. It appears at the observer-equation bridge-closure stage, where the explicit bridge correction tensor leaves the nonzero scalar projection \(\Pi_\phi^{(4)}\).
\end{corollary}

\begin{proof}
This is the direct combination of the positive same-class asymptotic closure and Theorem \ref{thm:verdict}.
\end{proof}

\section{Consequences for the Manuscript Program}

The present note changes the repository-wide decision tree again.

\begin{proposition}[Program verdict after explicit bridge-tensor reduction]
\label{prop:program}
After the explicit bridge-tensor reduction performed here, the surviving quartic branch should not be promoted to the active derivational line of \TheoryName{}.
\end{proposition}

\begin{proof}
Promotion to the active derivational line would require exact derivational bridge closure in the sense of the previous bridge-closure note. Theorem \ref{thm:verdict} shows that this requirement fails at current disciplined scope. Therefore the branch may remain recorded as a stable recovery-compatible candidate completion, but not as a successful derivation of the adopted Einsteinian observer equation.
\end{proof}

\begin{remark}
This is not a negative verdict on the same-class nonlinear-stability program. That program succeeded. It is a negative verdict on the stronger derivational claim. The distinction matters for the active manuscript sequence.
\end{remark}

\begin{corollary}[Next substantive burden if exact derivation remains the goal]
\label{cor:next}
If the active goal remains a true substrate derivation rather than only a stable recovery-compatible side program, the next substantive burden is now beyond the surviving quartic branch: one must either change the bridge/completion architecture or open a different candidate branch whose reduced observer equation can eliminate the residual tensor \(\mathcal C_{\mu\nu}^{\mathrm{br}}\) rather than merely suppress it.
\end{corollary}

\begin{proof}
By Proposition \ref{prop:program}, the present branch cannot be promoted to an exact derivational line. Therefore any continued derivational program must move beyond the current surviving quartic branch if it wants exact observer-equation closure.
\end{proof}

\section{What This Note Establishes}

This note establishes the following.
\begin{enumerate}
    \item It derives the reduced observer equation on the surviving quartic branch after auxiliary elimination.
    \item It computes the bridge correction tensor explicitly enough to identify it with the effective quartic scalar stress tensor.
    \item It gives the explicit reduced \(3+1\) form of that tensor.
    \item It computes the surviving flat-branch scalar projection \(\delta\mathcal C_\phi^{\mathrm{br}}=\Pi_\phi^{(4)}\phi\).
    \item It proves that the bridge correction tensor neither vanishes nor is pure gauge/constraint exact.
    \item It gives the derivational verdict: the surviving quartic branch is stable and recovery-compatible, but not exactly bridge-closing.
\end{enumerate}

It does not establish the following.
\begin{enumerate}
    \item It does not prove that no deeper completion anywhere in the \Aether{} / \Aflow{} program can ever derive Einsteinian gravity.
    \item It does not weaken the present ordered-flow positivity assumptions.
    \item It does not show that the surviving quartic branch is useless; it remains a mathematically disciplined stable candidate side program.
    \item It does not yet choose the next replacement branch or redesign architecture.
\end{enumerate}

\section{Conclusion}

The previous bridge-closure note asked the right exact question. After the same-class weighted/homogeneous route closes positively, what is the observer-equation correction tensor on the surviving quartic branch? Does it disappear? The present note now answers that question at the current disciplined scope.

The answer is no. After auxiliary elimination, the bridge correction tensor is the effective quartic scalar stress tensor. Its flat-branch scalar projection is the explicit nonzero self-energy \(\Pi_\phi^{(4)}\). That correction is small in the recorded recovery-compatible regime, but it is not zero and not pure gauge/constraint exact. Therefore the surviving quartic branch does not derive the exact Einsteinian observer equation, even though it remains stable and recovery-compatible.

This is the first genuine derivational obstruction on the current surviving branch. It appears only after the stability route succeeds. The manuscript program should therefore distinguish the two positive facts already earned from the negative derivational verdict now established: the branch is a disciplined stable candidate completion, but it is not the completed substrate derivation sought by the strongest version of the program.

\end{document}
